\subsection{\texorpdfstring{类型为 \(C_{l}\) 的代数 ( \(l \geq 1\) )}{类型为 C\_\{l\} 的代数 ( l \textbackslash geq 1 )}}

\begin{enumerate}
\def\labelenumi{(\Roman{enumi})}
\tightlist
\item
设 \(\varPsi\) 为定义在有限维向量空间 \(V\)（维数 \(2l \geq 2\)）上的非退化交错双线性型；\(x\) 的自同态 \(V\) 若满足 \(\varPsi(xv, v') + \varPsi(v, xv') = 0\)（对所有 \(v, v' \in V\)），则这些自同态构成 \(\mathfrak{sl}(V)\) 的半单李子代数，记为 \(\mathfrak{sp}(\varPsi)\)，称其为与 \(\varPsi\) 相关的辛李代数。
\end{enumerate}

由《代数》第 IX 章，§4，第2节，V 可写成两个极大全迷向子空间 F 和 \(F'\) 的直和，它们关于 \(\Psi\) 处于对偶关系。令 \((e_{i})_{1\leq i\leq l}\) 为 F 的一个基，令 \((e_{-i})_{1\leq i\leq l}\) 为 \(F'\) 的对偶基。于是

\[
(e _ {1}, \dots , e _ {l}, e _ {- l}, \dots , e _ {- 1})
\]

是 V 的一个基；我们称其为 V 的 Witt 基（或辛基）。相对于此基，\(\varPsi\) 的矩阵是阶为 2l 的方阵

\[
J = \left( \begin{array}{c c} 0 & s \\ - s & 0 \end{array} \right)
\]

其中 s 是阶为 l 的方阵，除第二对角线上的元素等于 1 外，其余元素全为 0，参见~第2.I节。

代数 \(\mathfrak{g} = \mathfrak{sp}(\varPsi)\) 可与方阵代数 \(\mathfrak{sp}(2l,k)\) 识别，后者由方阵 \(a\) 构成，方阵的阶为 \(2l\)，且满足 \(a = -J^{-1t}aJ = J^t aJ\)。

\[
a = \left( \begin{array}{c c} A & B \\ C & - s ^ {t} A s \end{array} \right)
\]

其中 A、B、C 是阶为 l 的方阵，并满足 \(B = s^{t}Bs\) 且 \(C = c^{t}Cs\)；换言之，B 和 C 关于第二对角线对称。因此

\[
\dim \mathfrak {g} = l ^ {2} + 2 \frac {l (l + 1)}{2} = l (2 l + 1).
\]

令 \(\mathfrak{h}\) 为 \(\mathfrak{g}\) 中对角矩阵的集合。这是 \(\mathfrak{g}\) 的交换子代数，其基由元素 \(H_{i} = E_{i,i} - E_{-i,-i}\)（\(1 \leq i \leq l\)）组成。令 \((\varepsilon_{i})_{1 \leq i \leq l}\) 为 \((H_{i})\) 的对偶基。对于 \(1 \leq i < j \leq l\)，置

\[
\left\{ \begin{array}{l l} X _ {2 \varepsilon_ {i}} & = E _ {i, - i} \\ X _ {- 2 \varepsilon_ {i}} & = - E _ {- i, i} \\ X _ {\varepsilon_ {i} - \varepsilon_ {j}} & = E _ {i, j} - E _ {- j, - i} \\ X _ {- \varepsilon_ {i} + \varepsilon_ {j}} & = - E _ {j, i} + E _ {- i, - j} \\ X _ {\varepsilon_ {i} + \varepsilon_ {j}} & = E _ {i, - j} + E _ {j, - i} \\ X _ {- \varepsilon_ {i} - \varepsilon_ {j}} & = - E _ {- i, j} - E _ {- j, i}. \end{array} \right.\tag{6}
\]

容易验证，这些元素构成 \(\mathfrak{h}\) 中 \(\mathfrak{g}\) 的一个补空间的基，并且对于 \(h\in \mathfrak{h}\)，有

\[
[ h, X _ {\alpha} ] = \alpha (h) X _ {\alpha}\tag{7}
\]

对所有 \(\alpha \in \mathbb{R}\) 都成立，其中 \(\mathbb{R}\) 是由 \(\pm 2\varepsilon_{i}\) 以及 \(\pm \varepsilon_{i} \pm \varepsilon_{j}\)（\(i < j\)）组成的集合。因此，\(\mathfrak{h}\) 等于其在 \(\mathfrak{g}\) 中的正规化子，因而是 \(\mathfrak{g}\) 的嘉当子代数；\(\mathfrak{h}\) 是分裂的，并且 \((\mathfrak{g}, \mathfrak{h})\) 的根正是 \(\mathbb{R}\) 的元素。\(\mathbb{R}\) 的根系 \((\mathfrak{g}, \mathfrak{h})\) 在 \(\mathrm{C}_{l}\) 时为 \(l \geq 2\) 型，在 \(\mathrm{A}_{1}\) 时为 \(\mathrm{C}_{1}\) 型（即 \(l = 1\) 型）（第 VI 章，§4，第6.I节，推广到 \(l = 1\) 的情形）。因此，\(\mathfrak{g}\) 是 \(\mathrm{C}_{l}\) 型的可分裂单李代数。

g 的每个分裂嘉当子代数都可由一个初等自同构变换为 h，因而也可由辛群 \(\mathbf{Sp}(\varPsi)\) 的某个元素变换得到（参见~(VII)）；所以它是集合 \(h_{\beta}\)，其中的元素 g 相对于 V 的某个 Witt 基 \(\beta\) 的矩阵为对角矩阵。立即验证可知，在 \(h_{\beta}\) 下不变的唯一向量子空间是由 \(\beta\) 的某个子集生成的子空间。

有 \(\mathfrak{sp}(2,k) = \mathfrak{sl}(2,k)\)。另一方面，代数 \(\mathfrak{sp}(4,k)\) 和 \(\mathfrak{o}_S(5,k)\) 具有相同的根系，因而同构（另参见~习题3）。下文设 \(l \geq 2\)。

\begin{enumerate}
\def\labelenumi{(\Roman{enumi})}
\setcounter{enumi}{1}
\tightlist
\item
根系 \(R^{\vee}\) 由第 VI 章，§4，第6.I节和第6.V节确定；我们得到
\end{enumerate}

\[
H _ {2 \varepsilon_ {i}} = H _ {i}, H _ {\varepsilon_ {i} - \varepsilon_ {j}} = H _ {i} - H _ {j}, H _ {\varepsilon_ {i} + \varepsilon_ {j}} = H _ {i} + H _ {j}.
\]

\begin{enumerate}
\def\labelenumi{(\Roman{enumi})}
\setcounter{enumi}{2}
\tightlist
\item
置 \(\alpha_{1} = \varepsilon_{1} - \varepsilon_{2},\ldots ,\alpha_{l - 1} = \varepsilon_{l - 1} - \varepsilon_{l},\alpha_{l} = 2\varepsilon_{l}\)。由第 VI 章，§4，第6.II节，\(\{\alpha_1,\dots ,\alpha_l\}\) 是 R 的一个基 B；相对于 B 的正根是 \(2\varepsilon_{i}\) 以及 \(\varepsilon_{i}\pm \varepsilon_{j}\)（\(i < j\)）。相应的 Borel 子代数 \(\mathfrak{b}\) 是 \(\mathfrak{g}\) 中上三角矩阵的集合。
\end{enumerate}

令 \(\delta\) 为 V 中的迷向旗（即其元素都是关于 \(\varPsi\) 的全迷向子空间），令 \(\mathfrak{p}_{\delta}\) 为由使 \(\mathfrak{g}\) 的各元素稳定的 \(\delta\) 中元素组成的子代数。像第2.III节中那样，我们证明映射 \(\delta \mapsto \mathfrak{p}_{\delta}\) 是从迷向旗的集合（分别从极大迷向旗的集合）到 \(\mathfrak{g}\) 的抛物子代数集合（分别到 Borel 子代数集合）的双射；当且仅当 \(\mathfrak{p}_{\delta} \supset \mathfrak{p}_{\delta'}\) 时，有 \(\delta \subset \delta'\)。

\begin{enumerate}
\def\labelenumi{(\Roman{enumi})}
\setcounter{enumi}{3}
\tightlist
\item
与 \(\alpha_{1},\ldots,\alpha_{l}\) 对应的基本权由第 VI 章，§4，第6.VI节给出，即 \(\varpi_{i}=\varepsilon_{1}+\cdots+\varepsilon_{i}\ (1\leq i\leq l)\)。
\end{enumerate}

我们将说明，权为 \(\sigma_{r}\) 的基本表示 \(\varpi_{r}\) 可以实现为 \(\bigwedge^{r}\sigma\) 的一个子表示，其中 \(\sigma\) 是 g 在 V 上的恒等表示；为此，我们将研究 g 在外代数 \(\bigwedge\sigma\) 上的表示 \(\bigwedge V\) 的分解。

令 \((e_{i}^{*})\) 为 \(V^{*}\) 中与 \((e_{i})\) 对偶的基。交错双线性型 \(\varPsi\) 可识别为 \(\varGamma^{*}\in\bigwedge^{2}V^{*}\) 中的元素，容易验证

\[
\Gamma^ {*} = - \sum_ {i = 1} ^ {l} e _ {i} ^ {*} \wedge e _ {- i} ^ {*}.
\]

令 \(\varPsi^{*}\) 为 \(\varPsi\) 的逆型（《代数》第 IX 章，§1，第7节）；立即有

\[
\varPsi^ {*} (e _ {i} ^ {*}, e _ {j} ^ {*}) = 0
\]

当 \(i \neq -j\) 时，\(\Psi^{*}(e_{i}^{*}, e_{-i}^{*}) = -1\)；且对 \(1 \leq i \leq l\)，\(\Psi^{*}\)。若将 \(\Gamma \in \bigwedge^{2}\mathrm{V}\) 视为相应的元素，则

\[
\Gamma = \sum_ {i = 1} ^ {l} e _ {i} \wedge e _ {- i}.
\]

记 \(X_{-}\) 为 \(\bigwedge V\) 上由以 \(\Gamma\) 作左外积给出的自同态，记 \(X_{+}\) 为 \(\bigwedge V\) 上由以 \(-\Gamma^{*}\) 作左内积给出的自同态：

\[
X _ {-} u = \left(\sum_ {i = 1} ^ {l} e _ {i} \wedge e _ {- i}\right) \wedge u,
\]

\[
X _ {+} u = \left(\sum_ {i = 1} ^ {l} e _ {i} ^ {*} \wedge e _ {- i} ^ {*}\right) \lrcorner u.
\]

为计算 \(X_{+}\) 和 \(X_{-}\)，以下述方式引入 \(\bigwedge V\) 的一个基：对于由 \([1, l]\) 的三个不相交子集组成的任意三元组 (A, B, C)，置

\[
e _ {\mathrm{A,B,C}} = e _ {a _ {1}} \wedge \dots \wedge e _ {a _ {m}} \wedge e _ {- b _ {1}} \wedge \dots \wedge e _ {- b _ {n}} \wedge e _ {c _ {1}} \wedge e _ {- c _ {1}} \wedge \dots \wedge e _ {c _ {p}} \wedge e _ {- c _ {p}}
\]

其中 \((a_{1},\ldots,a_{m})\)（分别为 \((b_{1},\ldots,b_{n}),(c_{1},\ldots,c_{p})\)）是 A（分别是 B、C）的元素按递增顺序排列所得的序列。这样得到 \(\bigwedge V\) 的一个基，简单计算表明

\[
X _ {-}. e _ {\mathrm{A,B,C}} = \sum_ {j \in \{1, l \}, j \notin \mathrm{A} \cup \mathrm{B} \cup \mathrm{C}} e _ {\mathrm{A,B,C} \cup \{j \}}\tag{8}
\]

\[
X _ {+}. e _ {\mathrm{A,B,C}} = - \sum_ {j \in \mathrm{C}} e _ {\mathrm{A,B,C-} \{j \}}.\tag{9}
\]

令 H 为 \(\bigwedge V\) 的自同态，它在 \((l - r)\) 上化为乘以 \(\bigwedge^{r}V\) 的乘法（\(0 \leq r \leq 2l\)）。容易验证（参见~习题19），有

\[
\begin{array}{c} {[ X _ {+}, X _ {-} ] = - H} \\ {[ H, X _ {+} ] = 2 X _ {+}} \\ {[ H, X _ {-} ] = - 2 X _ {-}.} \end{array}
\]

换言之，向量子空间 \(\mathfrak{s}\) 由 \(X_{+}, X_{-}\) 和 \(H\) 生成，且是 \(\operatorname{End}(\bigwedge \mathrm{V})\) 的李子代数，与 \(\mathfrak{sl}(2, k)\) 同构，而 \(\bigwedge^{r} \mathrm{V}\) 是权为 \(l - r\) 的元素子空间。令 \(\mathrm{E}_{r}\) 为 \(\bigwedge^{r} \mathrm{V}\) 中的本原元组成的子空间，即 \(\mathrm{E}_{r} = (\bigwedge^{r} \mathrm{V}) \cap \mathrm{Ker} X_{+}\)。由第1节可知，当 \(r < l\) 时，\(X_{-}\) 在 \(\bigwedge^{r} \mathrm{V}\) 上的限制是单射，并且当 \(r \leq l\) 时，\(\bigwedge^{r} \mathrm{V}\) 分解为直和。

\[
\begin{array}{c} \bigwedge^ {r} \mathrm{V} = \mathrm{E} _ {r} \oplus X _ {-} (\mathrm{E} _ {r - 2}) \oplus X _ {-} ^ {2} (\mathrm{E} _ {r - 4}) \oplus \dots \\ = \mathrm{E} _ {r} \oplus X _ {-} (\bigwedge^ {r - 2} \mathrm{V}). \end{array}
\]

这尤其表明，对于 \(\dim \mathrm{E}_r = \binom{2l}{r} - \binom{2l}{r-2}\)，\(0 \leq r \leq l\)。

另一方面，\(\mathfrak{sp}(\varPsi)\) 的定义本身表明，\(\Gamma^{*}\) 被 \(\sigma\) 的对偶的第二外幂湮灭。同样，\(\Gamma\) 被 \(\bigwedge^{2}\sigma\) 湮灭。立即得到，\(X_{+}\) 和 \(X_{-}\)，从而 \(H\)，也与 \(\bigwedge\sigma(g)\) 所对应的自同态 \(g\in\mathfrak{g}\) 交换。因此，对于 \(\mathrm{E}_{r}\)，子空间 \(0\leq r\leq l\) 在 \(\bigwedge^{r}\sigma\) 下稳定；我们将证明，\(\bigwedge^{r}\sigma\) 在 \(\mathrm{E}_{r}\) 上的限制是权为 \(\sigma_{r}\) 的基本表示 \(\varpi_{r}\)（\(1\leq r\leq l\)）。

我们先注意，\(\bigwedge^{r}\sigma\) 相对于 \(\mathfrak{h}\) 的权是

\[
\varepsilon_ {i _ {1}} + \dots + \varepsilon_ {i _ {k}} - (\varepsilon_ {j _ {1}} + \dots + \varepsilon_ {j _ {r - k}}),
\]

其中 \(i_1, \ldots, i_k\)（分别为 \(j_1, \ldots, j_{r-k}\)）是 \([1, l]\) 中的互异元素；因此，\(\bigwedge^r \sigma\) 的最高权确实是

\[
\varpi_ {r} = \varepsilon_ {1} + \dots + \varepsilon_ {r}
\]

表示 \(\varpi_{r}\) 的最高权向量是 \(e_{1} \wedge \cdots \wedge e_{r} = e_{\{1,\ldots,r\},\varnothing,\varnothing}\)，且它属于 \(e_{1} \wedge \cdots \wedge e_{r} \in E_{r}\)。由于 \(\bigwedge^{r}\sigma\) 是不可约表示，\(E_{r}\) 是其最高权

若 \(s \in \mathbf{Sp}(\varPsi)\)，则 s 在 \(\bigwedge V\)（分别在 \(\bigwedge V^{*}\)）上的延拓固定 \(\varGamma\)（分别固定 \(\varGamma^{*}\)），因而与 \(X_{+}\) 和 \(X_{-}\) 交换，并使 \(E_{r}\) 稳定。因此，\(E_{r}\) 包含由 \(F_{r}\) 将 \(e_{1} \wedge \cdots \wedge e_{r}\) 变换所得的向量生成的子空间 \(\mathbf{Sp}(\varPsi)\)。Witt 定理表明，这些变换得到的正是非零可分解 r-向量，且其对应的 V 的向量子空间是全迷向子空间；我们称这样的 r-向量为迷向的。

引理 2. 对于 \(1 \leq r \leq l\)，令 \(\mathrm{F}_r\) 为由迷向 \(\bigwedge^r\mathrm{V}\)-向量生成的 \(r\) 的子空间。则

\[
\bigwedge^ {r} \mathrm{V} = \mathrm{F} _ {r} + X _ {-} (\bigwedge^ {r - 2} \mathrm{V}) = \mathrm{F} _ {r} + \left(\sum_ {i = 1} ^ {l} e _ {i} \wedge e _ {- i}\right) \wedge \bigwedge^ {r - 2} \mathrm{V}.
\]

由于 \(F_{r} \subset E_{r}\)，且 \(\mathrm{E}_{r} \cap X_{-}(\bigwedge^{r-2}\mathrm{V}) = \{0\}\)，引理推出 \(F_{r} = E_{r}\)。另一方面，令 \(s \in \mathbf{Sp}(\varPsi)\)；自同构 \(a \mapsto sas^{-1}\) 将 \(\operatorname{End}(V)\) 保持为自身，保持 g 不变，并在其上诱导 \(\operatorname{Aut}_{0}(\mathfrak{g})\) 的一个元素（参见~(VII)），因而将 g 的每个不可约表示变为等价表示（第7节，第1节，命题2）。由于 \(e_{1} \wedge \cdots \wedge e_{r}\) 属于 \(\bigwedge^{r}\sigma\) 的一个不可约分支，且 \(E_{r} = F_{r}\) 由 \(e_{1} \wedge \cdots \wedge e_{r}\) 对 \(\mathbf{Sp}(\varPsi)\) 的变换生成，所以 g 在 \(E_{r}\) 上的表示是同构型的。但是，其最高权 \(\varpi_{r}\) 的重数为 1，故它不可约。

引理尚待证明。r = 1 时显然。作归纳并设 \(r \geq 2\)。根据归纳假设，只需证明

\[
\mathrm{F} _ {r - 1} \wedge \mathrm{V} \subset \mathrm{F} _ {r} + \Gamma \wedge \bigwedge^ {r - 2} \mathrm{V},
\]

或者说，若 \(y\) 是一个可分解的 \((r - 1)\) -向量且 \(x \in V\)，则

\[
z = y \wedge x \in \mathrm{F} _ {r} + \Gamma \wedge \bigwedge^ {r - 2} \mathrm{V}.
\]

令 \((f_{i})_{1\leq\pm i\leq l}\) 为 V 的 Witt 基，使 \(y=f_{1}\wedge\cdots\wedge f_{r-1}\)。只需在 \(x=f_{i}\) 时完成证明。若 \(i\notin(1-r,-1)\)，r-向量 \(f_{1}\wedge\cdots\wedge f_{r-1}\wedge f_{i}\) 是迷向的。否则，必要时重新编号 \(f_{i}\)，可设 i=1-r。于是 \(\Gamma=\sum_{j=1}^{l}f_{j}\wedge f_{-j}\)，所以

\[
\begin{array}{l} f _ {r - 1} \wedge f _ {1 - r} = \frac {1}{l - r + 2} \left(\Gamma - \sum_ {i = 1} ^ {r - 2} f _ {i} \wedge f _ {- i} + \sum_ {j = r} ^ {l} (f _ {r - 1} \wedge f _ {1 - r} - f _ {j} \wedge f _ {- j})\right), \\ z = \frac {1}{l - r + 2} \Gamma \wedge f _ {1} \wedge \dots \wedge f _ {r - 2} \\ \qquad + \frac {1}{l - r + 2} \sum_ {j = r} ^ {l} (f _ {1} \wedge \dots \wedge f _ {r - 2}) \wedge (f _ {r - 1} \wedge f _ {1 - r} - f _ {j} \wedge f _ {- j}). \end{array}
\]

但是

\[
f _ {r - 1} \wedge f _ {1 - r} - f _ {j} \wedge f _ {- j} = (f _ {r - 1} + f _ {j}) \wedge (f _ {1 - r} - f _ {- j}) - f _ {j} \wedge f _ {1 - r} + f _ {r - 1} \wedge f _ {- j}
\]

并立即验证下列 r-向量

\[
\begin{array}{c} f _ {1} \wedge \dots \wedge f _ {r - 2} \wedge (f _ {r - 1} + f _ {j}) \wedge (f _ {1 - r} - f _ {- j}), \\ f _ {1} \wedge \dots \wedge f _ {r - 2} \wedge f _ {j} \wedge f _ {1 - r} \text {and} f _ {1} \wedge \dots \wedge f _ {r - 1} \wedge f _ {- j} \end{array}
\]

当 \(r \leq j \leq l\) 时是迷向的。因此，\(z \in F_{r} + \Gamma \wedge \bigwedge^{r-2} V\)，从而证明完毕。

\begin{enumerate}
\def\labelenumi{(\Alph{enumi})}
\setcounter{enumi}{21}
\tightlist
\item
有 \(w_{0} = -1\)，所以 g 的每个有限维单表示都是正交型或辛型。由第 VI 章，§4，第6.VI节，\(\varpi_{r}\) 关于 \((\alpha_{1}, \ldots, \alpha_{l})\) 的坐标之和为
\end{enumerate}

\[
1 + 2 + \dots + (r - 1) + r + r + \dots + r + \frac {r}{2}
\]

因此 \(\sigma_{r}\) 在 r 为偶数时是正交型的，在 r 为奇数时是辛型的。

由于 \(e_{1} \wedge \cdots \wedge e_{r}\) 和 \(e_{-1} \wedge \cdots \wedge e_{-r}\) 属于 \(E_{r}\)，且

\[
\Psi_ {(r)} \left(e _ {1} \wedge \dots \wedge e _ {r}, e _ {- 1} \wedge \dots \wedge e _ {- r}\right) = 1,
\]

可知 \(\Psi_{(r)}\) 在 \(E_{r}\) 上的限制非零：这在相差一个非零常数因子的意义下就是该双线性型；当 r 为偶数时它对称，当 r 为奇数时它交错，并且在 \(\sigma_{r}\) 下不变。

\begin{enumerate}
\def\labelenumi{(\Roman{enumi})}
\setcounter{enumi}{5}
\tightlist
\item
对于所有 \(x \in \mathfrak{g}\)，\(\sigma(x)\) 的特征多项式具有形式
\end{enumerate}

\[
\mathrm{T} ^ {2 l} + f _ {1} (x) \mathrm{T} ^ {2 l - 1} + \dots + f _ {2 l} (x)
\]

其中 \(f_{1},\ldots ,f_{2l}\) 是 \(\mathfrak{g}\) 上的不变多项式函数。

若 \(x = \xi_{1}H_{1} + \cdots + \xi_{l}H_{l} \in \mathfrak{h}\)，则 \(f_{i}(x)\) 在相差一个符号的意义下是 \(\xi_{1}, \ldots, \xi_{l}, -\xi_{1}, \ldots, -\xi_{l}\) 的初等对称函数；这些对称函数在奇数次数时为零，并且

\[
\mathrm{T} ^ {2 l} + f _ {2} (x) \mathrm{T} ^ {2 l - 2} + \dots + f _ {2 l} (x) = (\mathrm{T} ^ {2} - \xi_ {1} ^ {2}) \ldots (\mathrm{T} ^ {2} - \xi_ {l} ^ {2}).
\]

像第2.VI节中那样，可知 \(f_{1}=f_{3}=f_{5}=\cdots=0\)，并且

\[
(f _ {2}, f _ {4}, \dots , f _ {2 l})
\]

是生成 \(\mathfrak{g}\) 上不变多项式函数代数的代数自由族。

\begin{enumerate}
\def\labelenumi{(\Roman{enumi})}
\setcounter{enumi}{6}
\tightlist
\item
由于邓金图的唯一自同构是恒等自同构，有 \(\mathrm{Aut}(\mathfrak{g}) = \mathrm{Aut}_0(\mathfrak{g})\)。
\end{enumerate}

令 \(\Sigma\) 为 \(V\) 关于 \(\Psi\) 的相似变换群。如第 2.VII 节中那样可证明，\(\mathfrak{g}\) 的自同构是映射 \(x \mapsto sxs^{-1}\)，其中 \(s \in \Sigma\)，所以 \(\operatorname{Aut}(\mathfrak{g}) = \operatorname{Aut}_0(\mathfrak{g})\) 可与 \(\Sigma / k^*\) 识别。

对于所有 \(s \in \Sigma\)，令 \(\mu(s)\) 为 s 的乘子。映射 \(s \mapsto \mu(s) \mod k^{*2}\) 从 \(\Sigma\) 到 \(k^{*}/k^{*2}\) 是同态，其核包含 \(k^{*}.1\)，因而给出从 \(\lambda\) 到 \(\Sigma/k^{*}\) 的同态 \(k^{*}/k^{*2}\)。有 \(\mathbf{Sp}(\varPsi) \cap k^{*} = \{1, -1\}\)。考虑同态序列

\[
1 \longrightarrow \mathbf {S p} (\Psi) / \{1, - 1 \} \stackrel {{\iota}} {{\longrightarrow}} \Sigma / k ^ {*} \stackrel {{\lambda}} {{\longrightarrow}} k ^ {*} / k ^ {* 2} \longrightarrow 1.\tag{10}
\]

映射 \(\iota\) 是单射，且 \(\mathrm{Im}(\iota) \subset \mathrm{Ker}\lambda\) 包含于 \(\mathbf{Sp}(\varPsi)\)，其元素的乘子为 1。如果 \(s \in \varSigma\) 的乘子是 \(k^{*2}\) 中的元素，则存在 \(\nu \in k^*\)，使 \(\nu s \in \mathbf{Sp}(\varPsi)\)；因此 \(\mathrm{Im}(\iota) = \mathrm{Ker}(\lambda)\)。总之，序列（10）是正合的。将 \(\mathbf{Sp}(\varPsi)/\{1, -1\}\) 识别为 \(\varSigma/k^*\) 的一个子群。由于 \(k^*/k^{*2}\) 是交换群，\(\mathbf{Sp}(\varPsi)/\{1, -1\}\) 包含 \(\varSigma/k^*\) 的换位子群。因此，\(\mathrm{Aut}_e(\mathfrak{g})\) 包含于 \(\mathbf{Sp}(\varPsi)/\{1, -1\}\)（\(\S\) 11，第2节，命题3）。事实上，它等于该群，并且 \(\mathrm{Aut}(\mathfrak{g})/\mathrm{Aut}_e(\mathfrak{g})\) 可与 \(k^*/k^{*2}\) 识别（习题9）。

\begin{enumerate}
\def\labelenumi{(\Roman{enumi})}
\setcounter{enumi}{7}
\tightlist
\item
典则双线性型 \(\Phi_{\mathrm{R}}\) 在 \(\mathfrak{h}^*\) 上为
\end{enumerate}

\[
\varPhi_ {\mathrm{R}} (\xi_ {1} \varepsilon_ {1} + \dots + \xi_ {l} \varepsilon_ {l}, \xi_ {1} ^ {\prime} \varepsilon_ {1} + \dots + \xi_ {l} ^ {\prime} \varepsilon_ {l}) = \frac {1}{4 (l + 1)} (\xi_ {1} \xi_ {1} ^ {\prime} + \dots + \xi_ {l} \xi_ {l} ^ {\prime})
\]

（第 VI 章，§4，第6.V节）。因此，\(\varPhi_{\mathrm{R}}\) 的逆型，即 Killing 型在 \(\mathfrak{h}\) 上的限制，为

\[
\Phi (\xi_ {1} H _ {1} + \dots + \xi_ {l} H _ {l}, \xi_ {1} ^ {\prime} H _ {1} + \dots + \xi_ {l} ^ {\prime} H _ {l}) = 4 (l + 1) (\xi_ {1} \xi_ {1} ^ {\prime} + \dots + \xi_ {l} \xi_ {l} ^ {\prime}).
\]

\begin{enumerate}
\def\labelenumi{(\Roman{enumi})}
\setcounter{enumi}{8}
\tightlist
\item
回顾由公式（6）定义的 \(X_{\alpha}\)（\(\alpha \in R\)）。容易验证，对于 \([X_{\alpha}, X_{-\alpha}] = -H_{\alpha}\)，有 \(\alpha \in R\)。另一方面，映射 \(\theta : a \mapsto -^{t}a\) 是 g 的自同构，且对所有 \(\theta(X_{\alpha}) = X_{-\alpha}\) 有 \(\alpha \in R\)。因此，\((X_{\alpha})_{\alpha \in \mathbb{R}}\) 是 \((\mathfrak{g}, \mathfrak{h})\) 中的谢瓦莱系。
\end{enumerate}

假设 \(k = \mathbf{Q}\)。嘉当子代数 \(\mathfrak{h}\) 有两个许可格：\(\mathrm{Q}(\mathrm{R}^{\vee}) = \sum_{i=1}^{l} \mathbf{Z}.H_i\) 和 \(\mathrm{P}(\mathrm{R}^{\vee}) = \mathrm{Q}(\mathrm{R}^{\vee}) + \frac{1}{2}\mathbf{Z}. \sum_{i=1}^{l} H_i\)。可见，\(\mathrm{Q}(\mathrm{R}^{\vee})\) 是 \(\mathfrak{h}\) 中具有整数项的矩阵集合。因此，谢瓦莱序 \(\mathrm{Q}(\mathrm{R}^{\vee}) + \sum_{\alpha \in \mathrm{R}} \mathbf{Z}.X_{\alpha}\) 是 \(\mathfrak{sp}(2l,\mathbf{Z})\)（其中矩阵属于 \(\mathfrak{g}\)）。

考虑约化李代数 \(\mathfrak{sp}(\varPsi) + \mathbf{Q}.1\)。容易看出，其具有整数项的元素组成一个谢瓦莱序；该序平行于 \(\mathfrak{sp}(\varPsi)\) 投影到 \(\mathbf{Q}.1\) 上，所得正是谢瓦莱序 \(\mathrm{P}(\mathrm{R}^{\vee}) + \sum \mathbf{Z}.X_{\alpha}\)。

最后，\(X_{\alpha}^{2} = 0\) 对所有 \(\alpha \in \mathbb{R}\) 都有。由此，格 \(\overline{\mathcal{V}}\) 位于 \(\mathrm{V}\) 中，并由 Witt 基 \(e_i\) 生成；它是 \(\mathfrak{sp}(2l,\mathbf{Z})\) 的容许格。在 \(\mathrm{E}_r \cap \bigwedge^r \mathcal{V}\) 中的 \(\mathrm{E}_r\) 也如此。

最后，\(E_{r}\) 对谢瓦莱序

\[
\mathrm{P} (\mathrm{R} ^ {\vee}) + \sum \mathbf {Z}. X _ {\alpha}
\]

仅当 \(r\) 为偶数时才有；此时 \(\mathrm{E}_r\cap \bigwedge^r\mathcal{V}\) 就是这样的格。

